% ============================================================
% ABSTRACT (no citations; last sentence = key finding; <= 1,920 characters for arXiv)
% ============================================================
Racial and ethnic minority patients receive less treatment for pain than
non-Hispanic White patients, and large language models (LLMs) could reproduce
that disparity when asked to assess pain. Prior audits tested a few proprietary
models, stated the patient's race explicitly and rarely measured the models' own
response instability. In Study~1, four open-weight models assessed six
acute-pain vignettes in every combination of four races and two sexes, with
identity stated or cued only by a name (two names per group), and each identity
was compared with a White male reference on the same case and sampling seed. In
Study~2, \nConfigs{} configurations of \nModels{} proprietary and open-weight
models, with reasoning on and off and at two precisions, answered the same
vignettes and the 50 clinician-written Q-Pain cases, with the case as the unit
of analysis. No difference survived correction in Study~1, none of \nTests{}
comparisons survived the false discovery rate in Study~2, and every severity
difference was equivalent to zero within one point on a 0--10 scale. The few
differences that survived correction within a model favoured Black patients.
Models differed far more from one another, from the clinical reference answer
and across their own reasoning and precision settings than between patients: the
reference patient received the high dose in 6\% to 93\% of responses depending
on the model. Audits should use several names per group, pair comparisons on the
seed, report the noise floor and test each deployed configuration. Current LLMs
showed no detectable racial or sex disadvantage in pain assessment, but which
model and settings a patient meets matters more than who the patient is.
